\section{基準エンコーダの選定}\label{app:encoder}
文献\cite{noguchi2022superposition}の図5に当たる他者視点の取り出しは，視覚エンコーダ2が A-1 の視界から作る符号と，視覚エンコーダ1に A-2 の視界を入れたときの符号がそろっていること（自他一致度が高いこと）に支えられている．
本研究の予備的な試みでは，運動以外の入力を加えて視覚エンコーダも学習すると，自他一致度が下がり，他者視点の取り出しが成り立たなかった（表\ref{tab:alignment}）．

\begin{table}[tbp]
  \centering
  \small
  \caption{視覚エンコーダの扱いと自他一致度}
  \label{tab:alignment}
  % paper/tables/alignment.tex -- generated by analyze/paper/tables_paper.py; sources in paper/tables/README.md
\begin{tabular}{@{}p{0.25\linewidth}p{0.27\linewidth}p{0.16\linewidth}p{0.2\linewidth}@{}}
\toprule
視覚エンコーダの扱い & 入力と環境・データ & 自他一致度 & 備考 \\
\midrule
一から学習（基準エンコーダ） & 運動だけ，入れ替えなし，1000エピソード，200エポック & 0.80（0.799） & 本研究で凍結して用いた重み．図5に当たる取り出しのマージン $+0.091$ \\
基準エンコーダを凍結 & 入れ替えあり，価値の信号（probe-Q）を入力 & 0.73 & \\
価値の入力で一から学習し直し & 入れ替えなし，価値の信号（probe-Q）を入力 & 0.25 & \\
基準エンコーダから微調整 & 入れ替えあり，価値の信号（probe-Q）を入力 & 0.26〜0.32 & \\
一から学習 & 快・不快だけを入力（運動なし），入れ替えあり & 0.16〜0.30 & 6本．図5に当たる取り出しは成り立たない \\
一から学習（試し学習） & 運動だけ，入れ替えあり，1000エピソード，150エポック & 0.47〜0.48 & \\
一から学習（試し学習） & 運動・快・不快・隠れフラグ，入れ替えあり，1000エピソード，150エポック & 0.18〜0.27 & 本文の「一から学習」の値 \\
一から学習（本番の規模） & 運動だけ，入れ替えあり，3000エピソード，200エポック & 0.740（ep100 で 0.783） & マージン $+0.075$．配置の区別は弱い \\
一から学習（本番の規模） & 運動・快・不快・隠れフラグ，入れ替えあり，3000エピソード，200エポック & 0.243 & マージン $+0.026$ \\
\bottomrule
\end{tabular}

  % 表A3：自他一致度（視覚エンコーダ2と視覚エンコーダ1の符号のコサイン）。
  % 行：基準エンコーダ（運動を入力，入れ替えなし，再現学習）0.80／基準エンコーダを凍結して入れ替えのある環境で学習 0.73／
  %   価値（外部の評価器の値）を入力してエンコーダも一から学習 0.25／入れ替えのある環境でエンコーダを微調整 0.26〜0.32／
  %   快・不快を入力して一から学習 0.16〜0.30／快・不快と運動の遮蔽を加えて一から学習（試し学習）0.18〜0.27／E1・E2（途中の値と明記）。
  % 出典：v11 まとめ §4 の 2（v4 0.80，v10 B3-frozen 0.73，v5 0.25，v10 B3-ft 0.26〜0.32，v11 0.16〜0.30），方法の節（0.18〜0.27）。
  % 版の番号は表に書かず，対応は tables/README.md に書く（第1弾の指示書 §4）。
\end{table}

自他一致度は，基準エンコーダで 0.80，文献\cite{noguchi2022superposition}の設定で学習したエンコーダを凍結して入れ替えのある環境で共有モジュールを学習した場合で 0.73 であった．
これに対し，運動以外の入力を加えてエンコーダも学習した場合は 0.16〜0.32 にとどまった．
入れ替えのない環境でも低下したので，入れ替えそのものが原因ではないと考えられる．
快・不快を入力として一から学習したモデルでは，他者視点の取り出しのマージンが $-0.012$〜$+0.013$ で，基準エンコーダ（$+0.091$）に遠く及ばなかった．
エンコーダの出力を線形に変換すると取り出しはほぼ回復したが，この変換は A-2 の本当の視界を使って学習したもので，情報があることの確認にとどまる．
% 出典：v11 まとめ §4 の 2，§5 の 8。

この比較は，入力・データ・学習の手順・損失が同時に違う複数のモデルを並べたものであり，運動以外の入力が原因であると操作して確かめたものではない．
本研究では，自他の対応を持ち，色の情報も弱いながら線形に読める唯一のエンコーダとして基準エンコーダを選び，すべての条件で凍結した．
自他の対応と色の表現を両立して学ぶエンコーダは得られておらず，本文5.4節の第二の限界に当たる．
% 出典：research_overview（v12 ログ §53 の判断 3，基準エンコーダの配置の情報は線形に読める程度）。
